\textbf{COLUNI 2025 -- Questão 16}\par\medskip Considere as expressões algébricas abaixo e as sentenças que se seguem:
\[
A=(a+b)^2;\qquad B=(a-b)^2;\qquad C=a^2+b^2;\qquad D=a^2-b^2
\]
\begin{enumerate}
\item[I.] Se $a=0$ e $b\in\mathbb{R}$, então $B=D$.
\item[II.] Se $a=\dfrac25$ e $b=\dfrac72$, então $A=15,21$.
\item[III.] Para todo $a,b\in\mathbb{R}$, têm-se que $A=C$.
\end{enumerate}
O número de sentenças CORRETAS é:
\begin{enumerate}
\item 0.
\item 1.
\item 2.
\item 3.
\end{enumerate}
